\chapter{Conclusion}\label{ch:conclusion}

This thesis demonstrated that a lightweight experimental platform is sufficient to teach core capacity-planning concepts using real measurements rather than purely theoretical examples. By subjecting a simple microservice to controlled workloads and collecting empirical traces, students can directly observe utilisation growth, queue buildup, saturation effects, and tail-latency amplification. The resulting measurements aligned closely with queueing-theoretic predictions and trace-driven discrete-event simulation (DES), confirming that the approach is both pedagogically effective and technically robust.

Beyond supporting learning outcomes in COMP9334, the work also addresses a clear gap in current capacity-planning practice and research. Existing tools either provide raw performance measurements without predictive modelling, or introduce modelling techniques without grounding them in real system behaviour. This thesis bridges that gap by demonstrating how empirical traces from a real containerised microservice can be combined with analytical methods and simulation. The framework establishes the foundations for a hybrid DES–ML modelling pipeline capable of predicting service degradation, tail-latency behaviour, and saturation thresholds under varying load and resource constraints—an emerging challenge in modern cloud and microservice environments.

The platform therefore serves a dual purpose. It provides students with an accessible, reproducible environment that runs on any personal computer, allowing them to explore complex performance phenomena in a controlled and cost-friendly setting. At the same time, it lays the groundwork for industrially relevant modelling techniques that can inform scaling decisions, resource-allocation policies, and performance forecasting in real systems. Extending the platform with ML-assisted service-time modelling and load-forecasting represents a promising direction for future work, further unifying educational objectives with contemporary capacity-engineering challenges.
